\documentclass[12pt]{article}
\usepackage{amsmath,amssymb,amsthm}
\usepackage[margin=1in]{geometry}

\title{Problem 164: Numbers for Which No Three Consecutive Digits\\Have a Sum Greater Than 9}
\author{Project Euler Solution}
\date{}

\begin{document}
\maketitle

\section{Problem Statement}
How many 20-digit numbers $n$ (without leading zeros) exist such that no three consecutive digits of $n$ have a digit sum greater than 9?

\section{Mathematical Analysis}

\subsection{Dynamic Programming Formulation}

Let $d_1 d_2 d_3 \ldots d_{20}$ be the digits of such a number. The constraint is:
\[
d_i + d_{i+1} + d_{i+2} \leq 9 \quad \text{for all } 1 \leq i \leq 18
\]
with $d_1 \geq 1$ (no leading zeros).

\subsection{State Definition}

Define $f(k, a, b)$ = number of valid completions for positions $k, k+1, \ldots, 20$, given that digit at position $k-2$ is $a$ and digit at position $k-1$ is $b$.

\subsection{Recurrence}

For position $k$ ($3 \leq k \leq 20$):
\[
f(k, a, b) = \sum_{d=0}^{9-a-b} f(k+1, b, d)
\]
where the upper limit ensures $a + b + d \leq 9$.

Base case: $f(21, a, b) = 1$ for all valid $(a, b)$.

\subsection{Answer Computation}

The first digit $d_1 \in \{1, \ldots, 9\}$, the second digit $d_2 \in \{0, \ldots, 9\}$:
\[
\text{Answer} = \sum_{d_1=1}^{9} \sum_{d_2=0}^{9} f(3, d_1, d_2)
\]

Note: when $k = 3$, the constraint $d_1 + d_2 + d_3 \leq 9$ applies, which is handled by the recurrence with $a = d_1, b = d_2$.

\subsection{State Space Analysis}

Valid states $(a, b)$ satisfy $0 \leq a, b \leq 9$. The constraint $a + b + d \leq 9$ requires $a + b \leq 9$ (otherwise no digit $d$ is available, meaning $f = 0$ for $a + b > 9$).

Number of valid states where $a + b \leq 9$:
\[
\sum_{a=0}^{9} (10 - a) = 55
\]

\subsection{Matrix Exponentiation (Alternative)}

The recurrence can be expressed as matrix multiplication on a 55-dimensional state vector. Let $\mathbf{v}_k$ be the vector indexed by valid pairs $(a,b)$ with $a+b \leq 9$, where the entry is $f(k, a, b)$.

The transition matrix $M$ has entry $M_{(a,b),(b,d)} = 1$ if $a + b + d \leq 9$ (and 0 otherwise, or if the indices don't chain properly).

Then $\mathbf{v}_k = M \cdot \mathbf{v}_{k+1}$, and the answer involves $M^{18} \cdot \mathbf{1}$ evaluated at appropriate initial states.

\section{Complexity}
\begin{itemize}
    \item \textbf{Time (iterative DP)}: $O(20 \times 10^3) = O(20000)$
    \item \textbf{Time (matrix exponentiation)}: $O(55^3 \log 20) \approx O(700000)$
    \item \textbf{Space}: $O(10^2) = O(100)$ with rolling array
\end{itemize}

\section{Answer}

\[
\boxed{378158756814587}
\]

\end{document}
